%! Author = Saeid Yazdani
%! Date = 5/18/2023

%\addcontentsline{toc}{section}{\bf Abstract (DE)}

\begin{flushleft}
%  \thispagestyle{empty}
    \section*{\textbf{\text{Kurzfassung}}}
    \justifying
    \vskip 1cm

    Elektrische und funktionale Tests sowie Qualifizierungen sind
    wesentliche Schritte im Produktionszyklus von Halbleiterbauelementen.
    Angesehene Hersteller setzen dabei stark auf \acs{ATE}, um während und
    nach der Fertigung vielfältige Tests durchzuführen und so die Lieferung
    sicherer und zuverlässiger Produkte zu gewährleisten.

    Mit der stetig wachsenden Komplexität von Halbleitern und den
    Fortschritten in Herstellungsprozessen steigt auch der Bedarf
    an hochentwickelten Test- und Messgeräten.
    Diese müssen in der Lage sein, sowohl sehr geringe
    Spannungen und Ströme bei integrierten Schaltkreisen als auch hohe
    Werte bei Leistungskomponenten mit höchster Präzision und Genauigkeit
    zu stimulieren und zu messen.

    Traditionelle \acs{ATE}-Systeme verlassen sich hauptsächlich auf analoge
    Regelkreise, um Testsignale innerhalb vorgegebener Grenzen zu halten.
    Allerdings weisen diese analogen Regelkreise bekannte Einschränkungen
    auf, die den modernen Anforderungen an Halbleitertests und
    Halbleiterqualifizierungen nicht immer gerecht werden.

    Zu diesen Einschränkungen gehören unter anderem ein hoher
    Bauteilaufwand, das Auftreten von Über- oder Unterschwingungen
    bei Laständerungen sowie der Bedarf an zusätzlicher Schaltungstechnik
    für Überwachungsfunktionen wie Überspannungs- oder Überstromschutz.

    Diese Arbeit analysiert die Grenzen herkömmlicher Regelkreise in Test-
    und Messgeräten und stellt ein hybrides Analog-Digital-Regelungskonzept
    vor, das die Schwächen rein analoger Systeme überwindet.

    Hierzu wurde ein Vier-Quadranten-\acs{DPS}-Prototyp mit
    \SI{\pm320}{\volt} entwickelt, der kontinuierlich Ströme bis zu
    \SI{2}{\ampere} und im Pulsbetrieb bis zu \SI{10}{\ampere} liefern kann.
    Der Prototyp wurde in praxisnahen Szenarien getestet, in denen stabile
    Stromversorgung für das \acs{DUT} sowie präzise und genaue Tests zur
    Charakterisierung und Verifizierung erforderlich sind.

    \vskip 0.45cm
    \section*{\textbf{Schlagwörter}}
    {\setstretch{1.0}
    Automatisierte Testsysteme, Digital, Analog, Regelkreis, Test,
    Qualifizierung, Sicherheit, Zuverlässigkeit, Verifikation, Messung,
    Halbleiter, Stromversorgung, Integrierter Schaltkreis
    \par}

\end{flushleft}
